\chapter{Giới thiệu}
\label{chap:introduction}

\section{Đặt vấn đề}

\subsection{Bối cảnh}

Quản lý tài chính cá nhân đòi hỏi người dùng ghi nhận giao dịch đều đặn, phân loại đúng và hiểu được ý nghĩa của lịch sử thu--chi. Kiến thức tài chính có liên quan trực tiếp đến khả năng lập kế hoạch và ra quyết định dài hạn \cite{lusardi2014}. Tuy nhiên, công cụ ghi chép truyền thống thường yêu cầu nhiều thao tác: chọn loại giao dịch, nhập số tiền, chọn danh mục, ví và thời gian. Khi nhập liệu trở thành gánh nặng, dữ liệu dễ bị thiếu hoặc sai danh mục; mọi báo cáo phía sau vì thế kém tin cậy.

Đầu vào tự nhiên như ``ăn phở sáng nay 45k bằng Momo'' giảm số bước thao tác, nhưng tạo ra bài toán kỹ thuật mới. Hệ thống phải hiểu đơn vị tiền Việt Nam, ngày tương đối, nhiều giao dịch trong một câu, tên ví gần đúng và ngữ cảnh phân loại. Ảnh hóa đơn và giọng nói còn bổ sung nhiễu OCR/STT. Nếu ghi thẳng kết quả suy đoán vào cơ sở dữ liệu, một lỗi nhỏ có thể làm sai số dư, ngân sách và toàn bộ insight về sau.

\subsection{Vấn đề về phân tích và diễn giải}

Biểu đồ và bộ lọc trả lời tốt những câu hỏi xác định như ``tháng này đã chi bao nhiêu''. Giá trị nghiên cứu của PERFIN không nằm ở việc bọc câu truy vấn SQL bằng chatbot. Hệ thống cần phát hiện các hiện tượng khó quan sát bằng một số đơn lẻ: xu hướng tăng dần, ngày chi bất thường, tốc độ cạn dòng tiền, khoản phí định kỳ hoặc tiến độ mục tiêu bị lệch.

Các hiện tượng trên phải được tính bởi giải thuật có thể kiểm thử. LLM có thể hiểu câu hỏi và diễn giải kết quả, nhưng không phải nguồn chân lý cho phép tính tài chính. Giao toàn bộ phân tích cho LLM làm kết quả khó tái lập và tăng nguy cơ xuất hiện số liệu không có trong dữ liệu gốc.

\subsection{Bài toán nghiên cứu}

Đề tài tập trung trả lời ba câu hỏi:

\begin{enumerate}
  \item Làm thế nào chuyển đầu vào văn bản, ảnh và giọng nói thành giao dịch có cấu trúc mà vẫn bảo vệ tính đúng đắn của dữ liệu?
  \item Những giải thuật xác định nào phù hợp để tạo insight có thể giải thích từ lịch sử tài chính cá nhân?
  \item LLM nên tham gia ở đâu để tăng tính tự nhiên và cá nhân hóa mà không chiếm quyền tính toán hoặc tự ý thay đổi dữ liệu?
\end{enumerate}

\widereportfigure{01-system-context.pdf}{Sơ đồ ngữ cảnh và phạm vi hệ thống PERFIN. Người dùng tương tác với ứng dụng; các nhà cung cấp LLM, OCR và STT chỉ là dịch vụ hỗ trợ bên ngoài.}{fig:system-context}

Hình~\ref{fig:system-context} định vị PERFIN là hệ thống quản lý dữ liệu và giải thuật. Ngân hàng, ví dùng chung, xác thực production và hạ tầng sẵn sàng cao không thuộc phạm vi niên luận cơ sở ngành.

\section{Mục tiêu}

Mục tiêu tổng quát là xây dựng và đánh giá một nguyên mẫu quản lý tài chính cá nhân trong đó dữ liệu có cấu trúc và giải thuật xác định là nền tảng, còn LLM là lớp hiểu ngôn ngữ và diễn giải có kiểm soát.

\begin{table}[H]
  \centering
  \caption{Mục tiêu và tiêu chí đánh giá của đề tài}
  \label{tab:objectives}
  \begin{tabularx}{\textwidth}{@{}L{1.1cm}Y Y@{}}
    \toprule
    \textbf{Mã} & \textbf{Mục tiêu cụ thể} & \textbf{Bằng chứng và tiêu chí nghiệm thu} \\
    \midrule
    O1 & Xây dựng mô hình dữ liệu tài chính có khóa, ràng buộc, chỉ mục và giao dịch nguyên tử. & Migration, ERD, test rollback; luồng cập nhật trọng yếu phải hoàn tất toàn bộ hoặc không đổi dữ liệu. \\
    O2 & Xây dựng pipeline trích xuất đa phương thức có hỏi lại và xác nhận. & Bộ dữ liệu gán nhãn, log tool call; không ghi DB khi thiếu trường bắt buộc. \\
    O3 & Hiện thực trend, anomaly, runway, recurring, correlation, ngân sách và mục tiêu. & Test công thức với dữ liệu thường, dữ liệu biên và kết quả tính tay. \\
    O4 & Giới hạn LLM ở hiểu ý định, gọi công cụ và diễn giải facts. & Tool schema, kiểm tra numeric faithfulness và fallback; không thêm số ngoài facts. \\
    O5 & Đánh giá khả năng vận hành của nguyên mẫu. & Test tích hợp, p50/p95, lỗi provider và UAT phải có log tái lập trước khi công bố. \\
    \bottomrule
  \end{tabularx}
\end{table}

Các ngưỡng trong Bảng~\ref{tab:objectives} là \textbf{mục tiêu nghiệm thu}, không mặc nhiên là kết quả đã đạt. Báo cáo chỉ gắn nhãn ``đã đo'' khi có lệnh, môi trường và đầu ra tái lập.

\section{Phạm vi đề tài}

\subsection{Phạm vi thực hiện}

\begin{table}[H]
  \centering
  \caption{Phạm vi thực hiện và nội dung ngoài phạm vi}
  \label{tab:scope}
  \begin{tabularx}{\textwidth}{@{}Y Y@{}}
    \toprule
    \textbf{Trong phạm vi} & \textbf{Ngoài phạm vi} \\
    \midrule
    Giao dịch thu, chi, chuyển ví và lãi/lỗ đầu tư ở mức nguyên mẫu. & Kết nối Open Banking và đồng bộ ngân hàng thật. \\
    Nhập văn bản, ảnh hóa đơn, giọng nói; preview và xác nhận. & Huấn luyện mô hình nền tảng, OCR hoặc STT mới. \\
    Danh mục, phản hồi sửa sai, gợi ý và re-tag có xác nhận. & Ví dùng chung, chia nợ nhóm và cộng tác thời gian thực. \\
    Ngân sách, recurring bill, insight và mục tiêu tài chính. & Tư vấn đầu tư, chấm điểm tín dụng hoặc quyết định thay người dùng. \\
    Một hồ sơ demo \code{default_user} và đường nâng cấp dữ liệu. & Đăng ký, JWT, phân quyền và cô lập multi-user production. \\
    Redis, worker và fallback cho demo/kiểm thử. & High availability, giám sát tập trung và cam kết uptime. \\
    \bottomrule
  \end{tabularx}
\end{table}

Trọng tâm của niên luận là mô hình dữ liệu, tính đúng của biến đổi dữ liệu và các giải thuật phân tích. Giao diện di động đóng vai trò thu nhận đầu vào và minh họa khả năng sử dụng; độ đầy đủ CRUD cho mọi đối tượng không phải thước đo đóng góp học thuật chính.

\subsection{Giả định và giới hạn}

\begin{itemize}
  \item Đơn vị tiền mặc định là VND; dữ liệu tiền dùng kiểu số thập phân trong PostgreSQL. Số do LLM sinh không được dùng làm nguồn chuẩn.
  \item Kết luận thống kê chỉ có ý nghĩa khi số quan sát đủ lớn. Hệ thống phải giảm mức tin cậy hoặc từ chối kết luận khi dữ liệu quá ít.
  \item Persona chỉ thay đổi cách diễn đạt, không thay đổi facts, phép tính, quyền truy cập hay điều kiện xác nhận.
  \item Nguyên mẫu không thay thế chuyên gia tài chính; gợi ý chỉ hỗ trợ tổ chức và hiểu dữ liệu cá nhân.
\end{itemize}

\section{Đóng góp của đề tài}

Đề tài hướng tới năm đóng góp có thể kiểm chứng:

\begin{enumerate}
  \item mô hình dữ liệu có ràng buộc và các luồng cập nhật nguyên tử;
  \item pipeline kết hợp LLM, parser cục bộ, fuzzy matching và con người trong vòng lặp;
  \item bộ giải thuật phân tích xác định, tách khỏi lớp sinh ngôn ngữ;
  \item cơ chế grounded narration, trong đó LLM chỉ diễn giải facts đã tính;
  \item giao thức đánh giá tách tính đúng thuật toán, chất lượng AI, hiệu năng và trải nghiệm nhập liệu.
\end{enumerate}

\section{Bố cục báo cáo}

Chương~\ref{chap:introduction} trình bày bài toán, mục tiêu và phạm vi. Chương~\ref{chap:theory} xây dựng cơ sở lý thuyết về dữ liệu tài chính, LLM, OCR/STT và giải thuật phân tích. Chương~\ref{chap:results} đặc tả yêu cầu, thiết kế kiến trúc, mô hình dữ liệu, luồng xử lý và kiểm thử. Chương~\ref{chap:conclusion} đối chiếu mục tiêu, nêu hạn chế và đề xuất hướng phát triển.
